\chapter*{Conclusion Générale et Perspectives}
\label{chap:conclusion}
\markboth{\MakeUppercase{Conclusion Générale et Perspectives}}{}
\addcontentsline{toc}{chapter}{Conclusion Générale et Perspectives}

Ce projet de fin d'études et de stage d'ingénieur a permis d'aborder une problématique industrielle et logicielle de tout premier plan : la digitalisation, l'automatisation et l'ordonnancement intelligent de la maintenance des équipements au sein d'un Progiciel de Gestion Intégré (ERP) d'entreprise, renforcé par les apports récents de l'intelligence artificielle générative et de l'automatisation des tests de navigation.

Au terme de ce travail mené au sein de la société suisse \textbf{\companyName} et de sa filiale d'ingénierie \textbf{\companySubsidiary} (Groupe \companyGroup), l'ensemble des objectifs fixés au cahier des charges ont été pleinement atteints. Le module \textbf{\odooModule} a été conçu, développé, refactorisé et validé selon les plus hauts standards d'ingénierie logicielle pour la plateforme de nouvelle génération \textbf{\odooVersion}.

\section*{Synthèse des Contributions et Résultats Clés}
Les réalisations majeures issues de ce projet s'articulent autour des axes suivants :

\begin{itemize}[leftmargin=1.5cm, label=\textcolor{odooPurple}{\small\textbullet}]
    \item \textbf{Une Solution ERP Industrielle Exhaustive et Modulaire} : Développement d'un système complet totalisant plus de \textbf{13\,700 lignes de code} réparties sur \textbf{76 fichiers}, intégrant \textbf{16 modèles de données ORM}, une hiérarchie spatiale multi-sites (\texttt{equipment.location}), la génération de codes QR pour identification mobile et des rapports vectoriels QWeb PDF.
    \item \textbf{Aide à la Décision \& Analytique Prédictive Réaliste} : Élaboration d'une formule multicritère de score de santé (\textit{Health Score} sur 100), couplée à des modèles statistiques de prévision des risques de panne (MTBF/MTTR), de projection budgétaire par moyennes mobiles pondérées et d'aide au remplacement technico-économique de machines.
    \item \textbf{Couche IA Conversationnelle BYO-LLM Sécurisée} : Intégration d'une passerelle multi-fournisseurs (Google Gemini et Cerebras) permettant la journalisation d'interventions en langage naturel et l'assistance interactive, sécurisée par une isolation stricte des outils et un flux obligatoire de confirmation explicite avant écriture.
    \item \textbf{Validation Complète de Navigation sous Playwright} : Déploiement d'une suite de \textbf{32 tests E2E} sous navigateur Chromium réel (16 fichiers de spécifications, 100\% de succès), ayant permis d'identifier et de corriger \textbf{10 anomalies réelles de navigation et d'interface} (dépliage des listes groupées Odoo 18, suppression de l'erreur 500 sur le Gantt, correction de visibilité des pièces de rechange, rechargement des fenêtres modales).
    \item \textbf{Télémétrie d'Assurance Qualité et Preuves sur Disque} : Validation sans faille attestée par \textbf{138 tests backend} (100\% de succès sur base fraîche \texttt{bt\_final3}), \textbf{46 tests mock LLM}, la validation live de l'API Gemini et l'enregistrement de deux vidéos complètes de démonstration.
\end{itemize}

\section*{Perspectives d'Évolution et Feuille de Route}
Afin de prolonger cette dynamique d'innovation technologique, plusieurs perspectives d'évolution sont envisagées :

\begin{enumerate}[leftmargin=1.2cm, label=\textbf{\arabic*.}]
    \item \textbf{Interconnexion IoT et Maintenance Conditionnelle (CBM)} : Interfacer le module avec des automates et capteurs industriels (via les protocoles MQTT et Modbus) pour ingérer en continu les données physiques (vibrations, température, ampérage, compteurs horaires) et déclencher automatiquement des ordres préventifs en cas de dérive anormale.
    \item \textbf{Entraînement de Modèles Supervisés de Machine Learning} : Exploiter les historiques d'interventions et les caractéristiques d'équipements extraites par le modèle \texttt{equipment.ai.data} pour entraîner des modèles de classification supervisée (\textit{Random Forest}, \textit{Gradient Boosting}) prédisant les modes de défaillance spécifiques.
    \item \textbf{Application Mobile Déconnectée pour Techniciens Nomades} : Concevoir une application mobile avec moteur de synchronisation bidirectionnelle hors-ligne pour les interventions sur des sites industriels ou des salles blanches dépourvues de connectivité réseau.
\end{enumerate}

En conclusion, ce stage a constitué une expérience professionnelle et scientifique particulièrement enrichissante au sein de l'écosystème Odoo suisse et tunisien. Il m'a offert l'opportunité de mobiliser les compétences acquises durant mon cursus d'ingénieur en génie logiciel et de délivrer un produit opérationnel, robuste, élégant et à très haute valeur ajoutée pour \textbf{\companyName}.